% Presentacion - PFM02 / Sesion 02
% Suma y resta con fracciones

\pdfminorversion=7
\documentclass[
	spanish,
	aspectratio=1610,
	xcolor={dvipsnames,table}
]{beamer}

\geometry{paperwidth=19.2cm,paperheight=12cm}

\usepackage[utf8]{inputenc}
\usepackage[T1]{fontenc}
\usepackage[spanish,es-nodecimaldot]{babel}
\usepackage{lmodern}
\usepackage{booktabs}
\usepackage{array}
\usepackage{graphicx}
\usepackage{ragged2e}
\usepackage{tikz}
\usetikzlibrary{arrows.meta,positioning,calc,decorations.pathreplacing,shapes.geometric,fit,backgrounds,overlay-beamer-styles}
\usepackage{hyperref}

% Datos del alumno y del programa
\newcommand{\studentname}{Mtro. Martín Jonathan de la Cruz Muñoz}
\newcommand{\studentshort}{M. de la Cruz Munoz}
\newcommand{\studentid}{00874}
\newcommand{\studentemail}{martin.delacruz@campus.iiiepe.edu.mx}
\newcommand{\programname}{Maestria en Aprendizaje y Ensenanza de las Matematicas}
\newcommand{\programshort}{MAEM}
\newcommand{\coursecode}{PFM02}
\newcommand{\coursename}{Temas selectos de matemáticas I}
\newcommand{\coursegroup}{Grupo A}
\newcommand{\courseperiod}{Marzo--julio 2026}
\newcommand{\institutionname}{IIIEPE}
\newcommand{\institutionlongname}{Instituto de Investigacion, Innovacion y Estudios de Posgrado para la Educacion del Estado de Nuevo Leon}
\newcommand{\institutionplace}{Monterrey, Nuevo Leon}
\newcommand{\iiiepelogo}{img/iiiepe/logo-iiiepe-blanco.png}
\newcommand{\iiiepemark}{img/iiiepe/marca-iiiepe.png}
\newcommand{\iiiepedots}{img/iiiepe/logo-puntos.png}

% Paleta institucional
\definecolor{iiiepeNavy}{HTML}{163B5C}
\definecolor{iiiepeBlue}{HTML}{1F6F9F}
\definecolor{iiiepeAqua}{HTML}{20B7A8}
\definecolor{iiiepeOrange}{HTML}{F28C28}
\definecolor{iiiepePaper}{HTML}{F6F8F7}
\definecolor{iiiepeInk}{HTML}{1E2A32}
\definecolor{iiiepeMuted}{HTML}{667085}

\mode<presentation>{
	\usetheme{default}
	\usefonttheme{professionalfonts}
	\setbeamertemplate{navigation symbols}{}
	\setbeamertemplate{blocks}[rounded][shadow=false]
}

\hypersetup{
	colorlinks=true,
	urlcolor=iiiepeBlue,
	linkcolor=iiiepeNavy
}

\setbeamercolor{background canvas}{bg=white}
\setbeamercolor{normal text}{fg=iiiepeInk,bg=white}
\setbeamercolor{structure}{fg=iiiepeBlue}
\setbeamercolor{frametitle}{fg=white,bg=iiiepeNavy}
\setbeamercolor{title}{fg=white,bg=iiiepeNavy}
\setbeamercolor{block title}{fg=white,bg=iiiepeBlue}
\setbeamercolor{block body}{fg=iiiepeInk,bg=iiiepePaper}
\setbeamercolor{alerted text}{fg=iiiepeOrange}
\setbeamerfont{title}{size=\LARGE,series=\bfseries}
\setbeamerfont{frametitle}{size=\fontsize{14}{18}\selectfont,series=\bfseries}
\setbeamerfont{block title}{series=\bfseries}
\setbeamertemplate{itemize item}{\textcolor{iiiepeAqua}{\fontsize{14}{18}\selectfont$\blacktriangleright$}}
\setbeamertemplate{itemize subitem}{\textcolor{iiiepeOrange}{\scriptsize$\blacksquare$}}
\setbeamertemplate{enumerate item}{\textcolor{iiiepeBlue}{\insertenumlabel.}}

\newcommand{\localLogo}[1]{%
	\IfFileExists{#1}{\includegraphics[height=0.58cm]{#1}}{\textbf{\institutionname}}%
}

\setbeamertemplate{frametitle}{
	\nointerlineskip
	\begin{beamercolorbox}[wd=\textwidth,ht=1.05cm,dp=0.22cm,leftskip=0.35cm,rightskip=0.25cm]{frametitle}
		\begin{minipage}[c]{0.72\textwidth}
			\insertframetitle
		\end{minipage}
		\hfill
		\raisebox{-0.1cm}{\localLogo{\iiiepelogo}}
	\end{beamercolorbox}
	\vspace{-0.04cm}
	{\color{iiiepeOrange}\rule{\textwidth}{1.4pt}}
}

\setbeamertemplate{footline}{
	\leavevmode
	\hbox{%
		\begin{beamercolorbox}[wd=.34\paperwidth,ht=2.7ex,dp=1ex,leftskip=1em]{author in head/foot}
			\color{white}\studentshort
		\end{beamercolorbox}
		\begin{beamercolorbox}[wd=.38\paperwidth,ht=2.7ex,dp=1ex,center]{title in head/foot}
			\color{white}\programshort\ -- \coursecode
		\end{beamercolorbox}
		\begin{beamercolorbox}[wd=.20\paperwidth,ht=2.7ex,dp=1ex,rightskip=1em plus 1fil]{date in head/foot}
			\color{white}Matricula \studentid\hfill\insertframenumber/\inserttotalframenumber
		\end{beamercolorbox}
	}
}
\setbeamercolor{author in head/foot}{bg=iiiepeNavy}
\setbeamercolor{title in head/foot}{bg=iiiepeBlue}
\setbeamercolor{date in head/foot}{bg=iiiepeNavy}

\newcommand{\accentbar}{%
	\begin{tikzpicture}[remember picture,overlay]
		\path[use as bounding box] (0,0) rectangle (0,0);
		\fill[iiiepeAqua] (current page.south west) rectangle ([xshift=0.42\paperwidth,yshift=0.08cm]current page.south west);
		\fill[iiiepeOrange] ([xshift=0.42\paperwidth]current page.south west) rectangle ([xshift=\paperwidth,yshift=0.08cm]current page.south west);
	\end{tikzpicture}
}

\newcommand{\DividirEntre}[2]{\ensuremath{#1\,\text{ entre }\,#2}}

\tikzset{
	cajaProc/.style={
		draw=iiiepeBlue,
		rounded corners=2mm,
		align=center,
		inner sep=4pt,
		minimum height=8mm,
		font=\small
	},
	cajaDato/.style={
		cajaProc,
		fill=iiiepeBlue!8
	},
	cajaOp/.style={
		cajaProc,
		fill=iiiepeOrange!12
	},
	cajaRes/.style={
		cajaProc,
		fill=iiiepeAqua!12
	},
	flechaProc/.style={
		-{Latex[length=2.4mm]},
		thick,
		draw=iiiepeBlue
	},
	notaProc/.style={
		draw=iiiepeOrange,
		rounded corners=2mm,
		fill=iiiepeOrange!12,
		align=center,
		inner sep=4pt,
		font=\scriptsize
	},
	decisionProc/.style={
		diamond,
		aspect=2.1,
		draw=iiiepeBlue,
		fill=iiiepeAqua!10,
		align=center,
		inner sep=2pt,
		font=\scriptsize
	}
}

\newcommand{\RutaOperacion}[4]{%
\begin{tikzpicture}[node distance=6mm]
	\node[cajaDato,visible on=<1->] (a) {#1};
	\node[cajaOp, right=of a,visible on=<2->] (b) {#2};
	\node[cajaOp, right=of b,visible on=<3->] (c) {#3};
	\node[cajaRes, right=of c,visible on=<4->] (d) {#4};
	\draw[flechaProc,visible on=<1->] (a) -- (b);
	\draw[flechaProc,visible on=<2->] (b) -- (c);
	\draw[flechaProc,visible on=<3->] (c) -- (d);
\end{tikzpicture}%
}

\newcommand{\sectionframe}[1]{%
	\begin{frame}[plain]
		\begin{tikzpicture}[remember picture,overlay]
			\fill[iiiepeNavy] (current page.south west) rectangle (current page.north east);
			\IfFileExists{\iiiepedots}{%
				\node[opacity=0.13,anchor=east] at ([xshift=0.8cm]current page.east) {\includegraphics[height=9.4cm]{\iiiepedots}};
			}{}
			\fill[iiiepeAqua] ([xshift=0.85cm,yshift=2.1cm]current page.south west) rectangle ([xshift=1.05cm,yshift=6.3cm]current page.south west);
			\fill[iiiepeOrange] ([xshift=1.15cm,yshift=2.1cm]current page.south west) rectangle ([xshift=1.35cm,yshift=5.35cm]current page.south west);
		\end{tikzpicture}
		\vspace{1.8cm}
		\hspace{1.25cm}
		\begin{minipage}{0.72\paperwidth}
			{\color{white}\Huge\bfseries #1}\\[0.25cm]
			{\color{white!75}Sesion 02 -- Suma y resta con fracciones}
		\end{minipage}
	\end{frame}
}

% Separadores automaticos desactivados para evitar diapositivas repetitivas.
% \AtBeginSection[]{\sectionframe{\insertsectionhead}}

\usepackage{tabularx,pgfplots,amssymb,adjustbox}\pgfplotsset{compat=1.18}\definecolor{iiiepePale}{HTML}{EEF5F8}
\title{Ecuaciones lineales y propiedades de la igualdad}\subtitle{Actividad 5}\begin{document}

\begin{frame}[plain]
\begin{tikzpicture}[remember picture,overlay]
\fill[iiiepeNavy] (current page.south west) rectangle ([xshift=.56\paperwidth]current page.north west);
\fill[iiiepeBlue] ([xshift=.56\paperwidth]current page.south west) rectangle (current page.north east);
\node[anchor=north west] at ([xshift=.55cm,yshift=-.45cm]current page.north west) {\localLogo{\iiiepelogo}};
\IfFileExists{\iiiepedots}{\node[opacity=.13,anchor=east] at (current page.east){\includegraphics[height=9cm]{\iiiepedots}};}{}
\end{tikzpicture}
\vspace{1.2cm}
\begin{minipage}{.50\paperwidth}
{\color{white}\fontsize{23}{27}\selectfont\bfseries\raggedright\inserttitle\par}\vspace{.35cm}
{\color{white!85}\fontsize{14}{18}\selectfont\insertsubtitle\par}\vspace{.5cm}
{\color{iiiepeOrange}\rule{.82\linewidth}{1.3pt}}\par\vspace{.4cm}
{\color{white}\studentname\par}\vspace{.1cm}
{\color{white!80}\small Matrícula: \studentid\par}
\end{minipage}\hfill
\begin{minipage}{.35\paperwidth}\raggedleft
{\color{white}\Large\bfseries IIIEPE\par}\vspace{.25cm}
{\color{white!85}\small\programname\par}\vspace{.5cm}
{\color{white}\small\coursename\par}\vspace{.5cm}
{\color{white!80}\small 13 de octubre de 2026\par Monterrey, Nuevo León\par}
\end{minipage}
\end{frame}
\begin{frame}{La igualdad como equilibrio}
\fontsize{14}{18}\selectfont\begin{itemize}\setlength{\itemsep}{.45cm}
\item 5 bolsas y 3 canicas equilibran 2 bolsas y 12 canicas.
\item 5x + 3 = 2x + 12; al retirar cantidades iguales, 3x = 9.
\item Cada bolsa contiene 3 canicas.
\end{itemize}
\end{frame}
\begin{frame}{Enteros y comprobación}
\fontsize{14}{18}\selectfont\begin{itemize}\setlength{\itemsep}{.45cm}
\item 5x + 3 = 2x + 12 $\rightarrow$ x = 3; ambos miembros valen 18.
\item 3x - 7 = 7x + 9 $\rightarrow$ x = -4; ambos miembros valen -19.
\item Aplicar la misma operación en ambos miembros conserva la igualdad.
\end{itemize}
\end{frame}
\begin{frame}{Paréntesis}
\fontsize{14}{18}\selectfont\begin{align*}4(2x+3)&=2(3x-5)\\8x+12&=6x-10\\2x&=-22\\x&=-11.\end{align*}
\centering Comprobación: $4(-22+3)=-76=2(-33-5)$.\\[.15cm]¿Qué operación preserva cada igualdad?
\end{frame}
\begin{frame}{Decimales y fracciones}
\fontsize{14}{18}\selectfont\begin{itemize}\setlength{\itemsep}{.45cm}
\item 0.5x + 2 = 0.2x + 3.8 $\rightarrow$ x = 6.
\item x/2 + 3 = x/4 + 5 $\rightarrow$ x = 8.
\item Multiplicar la segunda ecuación por 4 elimina sus denominadores.
\end{itemize}
\end{frame}
\begin{frame}{Apoyos ante rezago}
\fontsize{14}{18}\selectfont\begin{itemize}\setlength{\itemsep}{.45cm}
\item Usar balanza, fichas con signo y un ejemplo resuelto.
\item Cambiar una dificultad por vez, conservando el objetivo de equivalencia.
\item Permitir explicación oral antes de formalizar el procedimiento.
\end{itemize}
\end{frame}
\begin{frame}{Cierre formativo}
\fontsize{14}{18}\selectfont\begin{itemize}\setlength{\itemsep}{.45cm}
\item Devolver un logro, un error concreto y una acción de mejora.
\item Pedir corrección y comprobación, después un problema equivalente.
\item Recibir un archivo no demuestra que se comprendió la ecuación.
\end{itemize}
\end{frame}
\begin{frame}{Representación y explicación}
\centering\maxsizebox{.94\linewidth}{6.0cm}{\resizebox{.86\linewidth}{!}{\begin{tikzpicture}[font=\small,x=1cm,y=1cm]
\draw[thick] (-4,1)--(4,1);\draw[thick] (0,1)--(-.6,0)--(.6,0)--cycle;\draw[thick] (-3,1)--(-3,0);\draw[thick] (3,1)--(3,0);
\draw[fill=iiiepeBlue!15] (-4.3,-.7) rectangle (-1.7,0);\draw[fill=iiiepeAqua!20] (1.7,-.7) rectangle (4.3,0);
\node at (-3,-.35){$5x+3$};\node at (3,-.35){$2x+12$};\node at (0,-1.4){Restar $2x+3$ a ambos miembros: $3x=9$};\node at (0,-2){Dividir ambos miembros entre 3: $x=3$};\end{tikzpicture}}}\par\vspace{.15cm}\small Explica qué relación matemática muestra la representación.
\end{frame}
\begin{frame}{Reto de transferencia}
\fontsize{14}{18}\selectfont Compara $C_A=250+80h$ y $C_B=100+110h$: determina cuándo coinciden y explica tu decisión.\par\vspace{.5cm}Explica tu procedimiento, comprueba la respuesta y nombra una condición que podría cambiarla.
\end{frame}
\begin{frame}{Referencias}
\fontsize{9}{11}\selectfont\raggedright 
{\setlength{\parskip}{6pt}
\par\hangindent=1.27cm\hangafter=1  Academia Estatal de Matemáticas. (2023). \textit{Practicando saberes en Matemáticas: Cuaderno de trabajo. Primer grado. Secundaria}. Secretaría de Educación de Nuevo León.
\par\hangindent=1.27cm\hangafter=1  IIIEPE. (2026). \textit{Recurso sesión 5: Ecuaciones lineales y cuadráticas}. EVA posgrado.
\par\hangindent=1.27cm\hangafter=1  Secretaría de Educación Pública. (2024). \textit{Programa de estudio para la educación secundaria: Programa sintético de la Fase 6}. SEP.
}


\end{frame}
\end{document}
